\section{Actual curvature transport for closed two-colour pairs}
\label{sec:r9-closed-pair-curvature-transport}

The closedness of the entire PPD defect does not make its individual
summands closed.  The following finite homotopy identifies exactly
when a curvature factor can be moved to an extending reference form.
It also isolates a sufficient factorization of overlap defects that
would give an actual gluing witness.  That factorization is not
asserted for arbitrary analytic matrix paths.

\subsection{The pair complex and its finite actual homotopy}

Let $W$ be an analytic open over the original $K$.  For $s\ge1$ set
\[
 \begin{split}
 L_1^s(W)&=B^{s-1,s}(W)\oplus B^{s,s-1}(W),\\
 L_2^s(W)&=B^{s-2,s}(W)\oplus B^{s,s-2}(W)
                         \oplus B^{s-1,s-1}(W),\\
 \partial_1(p',p'')&=d'p'+d''p'',\\
 \partial_2(A,B,M)&=(d'A+d''M,\ d''B+d'M).
 \end{split}
 \tag{CPT1}
\]
Negative bidegree means zero.  All spaces are actual Mihatsch
spaces; the derivatives are the full derivatives.  The two-colour
signs give $\partial_1\partial_2=0$.  Write $Z_1^s(W)$ for
$\ker\partial_1$.  Objects are closed pairs, and equivalence is
difference by an actual displayed $\partial_2$ boundary.  One may
form the algebraic quotient $Z_1^s/\operatorname{im}\partial_2$,
but the formulas below construct the actual boundary witness and
never use a quotient element as a Green form.

Take an actual $h\in B^{k,k}(W)$, $k\ge0$, and
$p=(p',p'')\in Z_1^s(W)$.  Define the actual finite triple
\[
 H_h(p)=\bigl(d''h\wedge p',\ -d'h\wedge p'',\
                                  h\wedge d'p'\bigr)
                         \in L_2^{s+k+1}(W).
 \tag{CPT2}
\]
There is the literal actual identity
\[
                \partial_2 H_h(p)=(Dh\wedge p',\ Dh\wedge p'').
 \tag{CPT3}
\]
Indeed $h$ has even total degree, while $d'h,d''h$ have odd
degree.  The first component expands as
\[
 \begin{split}
 d'(d''h\wedge p')+d''(h\wedge d'p')
   &=Dh\wedge p'-d''h\wedge d'p'
                +d''h\wedge d'p'+h\wedge d''d'p'\\
   &=Dh\wedge p'-h\wedge Dp'.
 \end{split}
 \tag{CPT4}
\]
Closedness gives $d'p'=-d''p''$ and hence $Dp'=0$.
The second component is
\[
 d''(-d'h\wedge p'')+d'(h\wedge d'p')
  =Dh\wedge p''+d'h\wedge(d''p''+d'p')=Dh\wedge p''.
 \tag{CPT5}
\]
These are identities in the actual differential algebra, so no
test restriction, integration argument or current-to-form
injectivity is required.

All terms in (CPT2) have finite local actual presentations on
every common refinement.  Addition, actual restriction and
further refinement preserve (CPT3).  The pair spaces retain the
subspace topology from their own actual coefficient-final spaces,
and a quotient, if used, has its quotient topology.  No Hausdorff
or closed-image assertion about that quotient is made.  On a fixed
finite common coefficient presentation, with the support and
carriers fixed, the full Leibniz estimates control every mixed
coefficient derivative in (CPT2) by finitely many seminorms of
$h$ and $p$.  This is only a fixed-presentation operation statement;
it is not a global joint wedge-continuity theorem or a comparison
of a colimit image topology with an ambient inherited topology.

\subsection{Transport of closed curvature polynomials}

Suppose on $W$ that $R$ and $C_1,\ldots,C_a$ are actual
forms of degree $(1,1)$, closed under both full differentials,
and that actual scalar potentials satisfy
\[
                         C_j=R-D\Phi_j.
 \tag{CPT6}
\]
For a multi-index $\nu$ of total degree $d\ge1$, define
\[
 h_\nu=-\sum_{j:\nu_j>0}\Phi_j
   \left(\sum_{\ell=0}^{\nu_j-1}C_j^\ell R^{\nu_j-1-\ell}\right)
   \left(\prod_{i<j}C_i^{\nu_i}\right)
                            R^{\sum_{i>j}\nu_i}.
 \tag{CPT7}
\]
It is an actual form of degree $(d-1,d-1)$.  Every curvature
factor is closed, so telescoping gives
\[
                           Dh_\nu=C_1^{\nu_1}\cdots C_a^{\nu_a}-R^d.
 \tag{CPT8}
\]
For a homogeneous real polynomial $P=\sum_\nu t_\nu X^\nu$
of degree $d$, put $h_P=\sum_\nu t_\nu h_\nu$.  Then
\[
 Dh_P=P(C_1,\ldots,C_a)-P(1,\ldots,1)R^d.
 \tag{CPT9}
\]
For every closed pair $e\in Z_1^s(W)$, (CPT3) therefore gives
\[
 P(C_1,\ldots,C_a)e-P(1,\ldots,1)R^d e
                             =\partial_2 H_{h_P}(e).
 \tag{CPT10}
\]
Products in this formula multiply both components of the pair.
For degree zero the corresponding difference is zero.

Now let $W=N\setminus C$ and assume $R$ and $e$ are restrictions
of actual forms on all of $N$, with $e$ closed on $N$.
The second term in (CPT10) then actually extends over $C$.
Thus (CPT10) supplies an explicit boundary-plus-extending-pair
representation for the first term.  Neither the $\Phi_j$ nor
the $C_j$ need extend.  We do not claim that multiplication by
such a nonextending $C_j$ defines an endomorphism of a relative
quotient modulo extending pairs: that would require an extra
preservation assertion.  The result is the literal equation
(CPT10) for the specified inputs.

For the unit-symbol pair
\[
                 e_{a,b}=(a\,d''b,\ b\,d'a),
           \qquad Da=Db=0,
 \tag{CPT11}
\]
one has $\partial_1e_{a,b}=0$.  This follows from
$d'a\wedge d''b+d''b\wedge d'a=0$; the terms involving
$Da,Db$ vanish.  In particular $a=v(\lambda)$ and
$b=v(\mu)$ for original analytic units on $N$ give an
actual extending closed pair.  For $s\ge2$, setting
\[
 h=\Phi_0\sum_{\ell=0}^{s-2}C_0^\ell R^{s-2-\ell},
                \qquad Dh=R^{s-1}-C_0^{s-1},
 \tag{CPT12}
\]
gives the concrete transport
\[
       e_{a,b}R^{s-1}-e_{a,b}C_0^{s-1}=\partial_2H_h(e_{a,b}).
 \tag{CPT13}
\]
It identifies two candidate representatives of the known
commuting-square period.  It does not identify either of them
with the entire PPD defect modulo $\partial_2$.

\subsection{A checkable overlap consequence and the remaining gap}

Retain the actual CI data and open-cover setup of RAG.  Suppose
on a whole punctured overlap a proposed primitive difference
has the verified finite factorization
\[
         p_1-p_0=\partial_2q+
               \sum_{j=1}^t P_j(C_{j,1},\ldots,C_{j,a_j})e_j,
 \tag{CPT14}
\]
where each $e_j$ is an actual closed pair extending over the
unpunctured overlap, all bidegrees match, and each displayed
curvature admits a potential relative to one extending closed
reference $R$.  The equality in (CPT14) is an actual equality
on the entire overlap, not only a period identity or equality
in an unspecified quotient.  Then the explicit formulas above
produce
\[
 \begin{split}
 q_{\rm new}&=q+\sum_jH_{h_{P_j}}(e_j),\\
 e_{\rm new}&=\sum_jP_j(1,\ldots,1)R^{\deg P_j}e_j,\\
                       p_1-p_0&=\partial_2q_{\rm new}+e_{\rm new}|_W.
 \end{split}
 \tag{CPT15}
\]
RAG now applies with these actual witnesses and its full
partition-of-unity correction, followed by RSP to construct
the original-support actual Green form.  For more than two
opens, the RAG requirement of a verified witness on the whole
sequential overlap remains; pairwise witnesses are not silently
substituted for that data.

The conditional datum (CPT14) is finite and can be verified for
an explicitly supplied factorization independently of a Green
existence assertion.  It is not automatically present for an
arbitrary frame cover.  In PPD the complete defect is closed,
but the individual triangle, wall and metric-transport pairs
are not asserted separately closed.  Applying (CPT3) to them
term by term would retain the nonzero remainder $-hDp'$ in
(CPT4), or the second-component remainder in (CPT5).
Consequently this computation isolates the missing
factorization rather than removing it.  It proves no general
coherence, full lower-middle kernel comparison, actual descent
or Beilinson--Bloch height theorem.

One must also not cancel a curvature factor from
$\partial_1(C_jp)=0$ to conclude $\partial_1p=0$.
Actual wedge multiplication has dimension and support zero
divisors.  A regular-variable argument in a free polynomial
algebra does not supply such cancellation in the actual
quotient.  Closedness of each coefficient pair used in
(CPT10) or (CPT14) is an independent retained hypothesis.

For a direct calibration on a two-unit torus, let
$x=v(T_1)$, $y=v(T_2)$,
$C=D(x^2/2)=d'x\wedge d''x$, and $p=(x\,d''y,0)$.
Then $\partial_1p=d'x\wedge d''y\ne0$, whereas
$C\wedge\partial_1p=0$.  Thus $Cp$ is a closed actual
pair but $p$ is not.  This elementary wedge-zero-divisor
example is a warning about that inference, not a new
counterexample to strong G or to the frozen wedge statements.
